\documentclass[10pt,a4paper]{amsart}
\usepackage{amsmath,amssymb,amsthm,graphicx}
\usepackage{verbatim}
\newenvironment{reference}{\comment}{\endcomment}
%\newenvironment{reference}{}{}
\newenvironment{slogan}{\comment}{\endcomment}
\newenvironment{history}{\comment}{\endcomment}

% For commutative diagrams we use Xy-pic
\usepackage[all]{xy}

% We use 2cell for 2-commutative diagrams.
\xyoption{2cell}
% Theorem environments.
%
\theoremstyle{plain}
\newtheorem{theorem}[subsection]{Theorem}
\newtheorem{proposition}[subsection]{Proposition}
\newtheorem{lemma}[subsection]{Lemma}

\theoremstyle{definition}
\newtheorem{definition}[subsection]{Definition}
\newtheorem{example}[subsection]{Example}
\newtheorem{exercise}[subsection]{Exercise}
\newtheorem{situation}[subsection]{Situation}

\theoremstyle{remark}
\newtheorem{remark}[subsection]{Remark}
\newtheorem{remarks}[subsection]{Remarks}

\numberwithin{equation}{subsection}

% Macros
%
\def\lim{\mathop{\mathrm{lim}}\nolimits}
\def\colim{\mathop{\mathrm{colim}}\nolimits}
\def\Spec{\mathop{\mathrm{Spec}}}
\def\Hom{\mathop{\mathrm{Hom}}\nolimits}
\def\Ext{\mathop{\mathrm{Ext}}\nolimits}
\def\SheafHom{\mathop{\mathcal{H}\!\mathit{om}}\nolimits}
\def\SheafExt{\mathop{\mathcal{E}\!\mathit{xt}}\nolimits}
\def\Sch{\mathit{Sch}}
\def\Mor{\mathop{\mathrm{Mor}}\nolimits}
\def\Ob{\mathop{\mathrm{Ob}}\nolimits}
\def\Sh{\mathop{\mathit{Sh}}\nolimits}
\def\NL{\mathop{N\!L}\nolimits}
\def\CH{\mathop{\mathrm{CH}}\nolimits}
\def\proetale{{pro\text{-}\acute{e}tale}}
\def\etale{{\acute{e}tale}}
\def\QCoh{\mathit{QCoh}}
\def\Ker{\mathop{\mathrm{Ker}}}
\def\Im{\mathop{\mathrm{Im}}}
\def\Coker{\mathop{\mathrm{Coker}}}
\def\Coim{\mathop{\mathrm{Coim}}}

% Boxtimes
%
\DeclareMathSymbol{\boxtimes}{\mathbin}{AMSa}{"02}

%
% Macros for moduli stacks/spaces
%
\def\QCohstack{\mathcal{QC}\!\mathit{oh}}
\def\Cohstack{\mathcal{C}\!\mathit{oh}}
\def\Spacesstack{\mathcal{S}\!\mathit{paces}}
\def\Quotfunctor{\mathrm{Quot}}
\def\Hilbfunctor{\mathrm{Hilb}}
\def\Curvesstack{\mathcal{C}\!\mathit{urves}}
\def\Polarizedstack{\mathcal{P}\!\mathit{olarized}}
\def\Complexesstack{\mathcal{C}\!\mathit{omplexes}}
% \Pic is the operator that assigns to X its picard group, usage \Pic(X)
% \Picardstack_{X/B} denotes the Picard stack of X over B
% \Picardfunctor_{X/B} denotes the Picard functor of X over B
\def\Pic{\mathop{\mathrm{Pic}}\nolimits}
\def\Picardstack{\mathcal{P}\!\mathit{ic}}
\def\Picardfunctor{\mathrm{Pic}}
\def\Deformationcategory{\mathcal{D}\!\mathit{ef}}
\usepackage{cleveref}
\crefname{originallemma}{original lemma}{original lemmas}
\Crefname{originallemma}{Original lemma}{Original lemmas}
\makeatletter
\newlabel{algebra-lemma-tensor-localization}{{00DK}{source}{}{sourcepointer-1}{}}
\newlabel{algebra-lemma-tensor-localization@cref}{{[originallemma][0][]00DK}{[1][0][]source}{}{}{}}
\makeatother
\crefname{originalproposition}{original proposition}{original propositions}
\Crefname{originalproposition}{Original proposition}{Original propositions}
\makeatletter
\newlabel{algebra-proposition-localization-exact}{{00CS}{source}{}{sourcepointer-2}{}}
\newlabel{algebra-proposition-localization-exact@cref}{{[originalproposition][0][]00CS}{[1][0][]source}{}{}{}}
\makeatother
\crefname{originallemma}{original lemma}{original lemmas}
\Crefname{originallemma}{Original lemma}{Original lemmas}
\makeatletter
\newlabel{homology-lemma-quotient-by-kernel-exact-functor}{{06XK}{source}{}{sourcepointer-3}{}}
\newlabel{homology-lemma-quotient-by-kernel-exact-functor@cref}{{[originallemma][0][]06XK}{[1][0][]source}{}{}{}}
\makeatother
\crefname{originallemma}{original lemma}{original lemmas}
\Crefname{originallemma}{Original lemma}{Original lemmas}
\makeatletter
\newlabel{homology-lemma-serre-subcategory-is-kernel}{{02MS}{source}{}{sourcepointer-4}{}}
\newlabel{homology-lemma-serre-subcategory-is-kernel@cref}{{[originallemma][0][]02MS}{[1][0][]source}{}{}{}}
\makeatother
\title{Localization yields an equivalence Mod\_A / T -> Mod\_\{S\textasciicircum{}\{-1\}A\} for the Serre subcategory of S-torsion modules}
\author{The Stacks Project Authors}
\date{stacks-project:examples.tex:local-snapshot-20261012; GNU FDL 1.2 or later}
\begin{document}
\maketitle
\noindent\textit{Editorial note.} All delivered mathematical text is the author's own, copied from the pinned Stacks Project chapter file examples.tex; no statement or proof step was rewritten, reordered or removed.\par\medskip
\noindent\textit{Editorial note.} The counted proof is the author's complete argument: Fullness is witnessed by the same essential-surjectivity isomorphism; the author proves faithfulness and essential surjectivity and states the descended functor without a separate fullness claim - content verbatim.\par\medskip
\noindent\textit{Editorial note.} The label proposition-localization-and-serre-quotients names the counted statement, and the proof environment immediately following it in the source is delivered verbatim as its proof.\par\medskip
\noindent\textit{Editorial note.} The counted statement is carried by the author's own proposition environment in the chapter source.\par\medskip
\noindent\textit{Allowed standard prerequisites (editorial).}\begin{itemize}
\item Serre subcategories and quotient categories of abelian categories; the universal property of the quotient functor (Homology, serre-subcategory-is-kernel) and the faithful/quotient-by-kernel criterion (Homology, quotient-by-kernel-exact-functor)
\item Exactness of localization (Algebra, proposition-localization-exact) and S\textasciicircum{}\{-1\}M congruent to M tensor S\textasciicircum{}\{-1\}A (Algebra, tensor-localization)
\end{itemize}\medskip
\ifdefined\hypertarget\hypertarget{sourcepointer-1}{}\fi
\noindent\textit{Source reference.} Original lemma 00DK: sources/stacks-chow/tags.txt, source lines 500--500. Prerequisite pointer invoked at 6805; it names a standard fact of a sibling chapter that the counted argument uses, and every step of the counted proof built on it is retained verbatim in the delivered body.\par\medskip
\ifdefined\hypertarget\hypertarget{sourcepointer-2}{}\fi
\noindent\textit{Source reference.} Original proposition 00CS: sources/stacks-chow/tags.txt, source lines 472--472. Prerequisite pointer invoked at 6795; it names a standard fact of a sibling chapter that the counted argument uses, and every step of the counted proof built on it is retained verbatim in the delivered body.\par\medskip
\ifdefined\hypertarget\hypertarget{sourcepointer-3}{}\fi
\noindent\textit{Source reference.} Original lemma 06XK: sources/stacks-chow/tags.txt, source lines 8854--8854. Prerequisite pointer invoked at 6806; it names a standard fact of a sibling chapter that the counted argument uses, and every step of the counted proof built on it is retained verbatim in the delivered body.\par\medskip
\ifdefined\hypertarget\hypertarget{sourcepointer-4}{}\fi
\noindent\textit{Source reference.} Original lemma 02MS: sources/stacks-chow/tags.txt, source lines 3455--3455. Prerequisite pointer invoked at 6798; it names a standard fact of a sibling chapter that the counted argument uses, and every step of the counted proof built on it is retained verbatim in the delivered body.\par\medskip
\begin{proposition}
\label{proposition-localization-and-serre-quotients}
Let $A$ be a ring. Let $S$ be a multiplicative subset of $A$.
Let $\text{Mod}_A$ denote the category of $A$-modules and $\mathcal{T}$ its
Serre subcategory of modules for which any element is annihilated by some
element of $S$. Then there is a canonical equivalence
$\text{Mod}_A/\mathcal{T} \rightarrow \text{Mod}_{S^{-1}A}$.
\end{proposition}

\begin{proof}
The functor $\text{Mod}_A \to \text{Mod}_{S^{-1}A}$ given by $M
\mapsto M \otimes_A S^{-1}A$ is exact (by Algebra, Proposition
\ref{algebra-proposition-localization-exact})
and maps modules in $\mathcal{T}$ to zero.
Thus, by the universal property given in Homology, Lemma
\ref{homology-lemma-serre-subcategory-is-kernel}, the functor descends to a
functor $\text{Mod}_A/\mathcal{T} \to \text{Mod}_{S^{-1}A}$.

\medskip\noindent
Conversely, any $A$-module $M$ with $M \otimes_A S^{-1}A = 0$
is an object of $\mathcal{T}$, since
$M \otimes_A S^{-1}A \cong S^{-1} M$
(Algebra, Lemma \ref{algebra-lemma-tensor-localization}). Thus
Homology, Lemma \ref{homology-lemma-quotient-by-kernel-exact-functor}
shows that the functor
$\text{Mod}_A/\mathcal{T} \to \text{Mod}_{S^{-1}A}$ is faithful.

\medskip\noindent
Furthermore, this embedding is essentially surjective: a preimage to an
$S^{-1}A$-module $N$ is $N_A$, that is $N$ regarded as an $A$-module, since the
canonical map $N_A \otimes_A S^{-1}A \to N$ which maps $x \otimes a/s$ to
$(a/s) \cdot x$ is an isomorphism of $S^{-1}A$-modules.
\end{proof}

\end{document}
